% GENERATED FILE. Edit canonical lessons, then run npm run generate:books.
% canonical-source-hash: fnv1a64:dfddb9f4996bf5e2
% canonical-lessons: JA-C262-jiko, JA-C262-tasukaru, JA-C262-nigeru, JA-C262-totsuzen, JA-C262-higai

\chapter{An Accident by the Station}
\label{ch:ja-an-accident-by-the-station}
\hlchaptermodality{262}

\begin{chapteropening}
\textbf{By the end of this chapter:} \emph{I can say that there was an accident and who was saved.}

Complete the last lesson of chapter 262: I can say that there was an accident and who was saved.
\end{chapteropening}

\section[\texorpdfstring{\ja{じこ} (jiko)}{じこ (jiko)}]{\ja{じこ} (jiko) — an accident}

\label{lesson:JA-C262-jiko}



\begin{quote}
Before the new one: say the Japanese for serious, then the Japanese for selfish.
\end{quote}

\subsection*{You'll want to know: \ja{じこ}}
\textbf{\ja{じこ}} — \emph{jiko} — \textquotedblleft{}an accident\textquotedblright{}.

\begin{quote}
\textbf{\ja{きのう} \ja{えきの} \ja{ちかくで} \ja{じこが} \ja{ありました。}} — \emph{kinō eki no chikaku de jiko ga arimashita} — \textquotedblleft{}There was an accident near the station yesterday.\textquotedblright{}
\end{quote}

\begin{cousinweb}
\textbf{\ja{じこ}} (\emph{jiko}) is Sino-Japanese (kango): two Chinese characters read with their on'yomi, readings borrowed from Chinese. They are \emph{ji}, \textquotedblleft{}matter\textquotedblright{} (Middle Chinese \emph{dzriH}) and \emph{ko}, \textquotedblleft{}circumstances\textquotedblright{} (\emph{kuH}) (KANJIDIC2; Middle Chinese in Baxter's transcription).
\end{cousinweb}

\begin{grammarlens}[title={What you've built}]
One more word for telling what happened.
\end{grammarlens}

\subsection*{Your turn}
\begin{itemize}
\raggedright
  \item \emph{Say it:} \emph{jiko}
  \item \emph{Say it:} \emph{jiko}, once more
  \item \emph{Say it:} \emph{kinō eki no chikaku de jiko ga arimashita} — There was an accident near the station yesterday.
  \item \emph{Recall:} say the Japanese for serious, then the Japanese for selfish, then say \emph{jiko} again
\end{itemize}

\begin{tcolorbox}[breakable,colback=teal!4,colframe=teal!35!black,title={Before you move on}]
What does \ja{じこ} mean? (\textquotedblleft{}An accident\textquotedblright{}.) Say it once more.
\end{tcolorbox}
\section[\texorpdfstring{\ja{たすかる} (tasukaru)}{たすかる (tasukaru)}]{\ja{たすかる} (tasukaru) — to be saved; to survive}

\label{lesson:JA-C262-tasukaru}



\begin{quote}
Before the new one: say the Japanese for selfish, then the Japanese for an accident.
\end{quote}

\subsection*{You'll want to know: \ja{たすかる}}
\textbf{\ja{たすかる}} — \emph{tasukaru} — \textquotedblleft{}to be saved; to survive\textquotedblright{}.

\begin{quote}
\textbf{\ja{くるまの} \ja{じこで} \ja{ちちは} \ja{たすかりました。}} — \emph{kuruma no jiko de chichi wa tasukarimashita} — \textquotedblleft{}My father came through the car crash alive.\textquotedblright{}
\end{quote}

Polite \textbf{\ja{たすかります}}, \emph{tasukarimasu}; a godan verb. Past (polite): \textbf{\ja{こどもは} \ja{たすかりました。}}, \emph{kodomo wa tasukarimashita}.

\begin{cousinweb}
\textbf{\ja{たすかる}} (\emph{tasukaru}) is a native Japanese word (wago), inherited rather than borrowed from Chinese; its kanji was attached later as a kun'yomi (JMdict). It is the intransitive partner of \emph{tasukeru}, \textquotedblleft{}to help\textquotedblright{}.
\end{cousinweb}

\begin{grammarlens}[title={What you've built}]
One more word for telling what happened.
\end{grammarlens}

\subsection*{Your turn}
\begin{itemize}
\raggedright
  \item \emph{Say it:} \emph{tasukaru}
  \item \emph{Say it:} \emph{tasukaru}, once more
  \item \emph{Say it:} \emph{kuruma no jiko de chichi wa tasukarimashita} — My father came through the car crash alive.
  \item \emph{Recall:} say the Japanese for selfish, then the Japanese for an accident, then say \emph{tasukaru} again
\end{itemize}

\begin{tcolorbox}[breakable,colback=teal!4,colframe=teal!35!black,title={Before you move on}]
What does \ja{たすかる} mean? (\textquotedblleft{}To be saved; to survive\textquotedblright{}.) Say it once more.
\end{tcolorbox}
\section[\texorpdfstring{\ja{にげる} (nigeru)}{にげる (nigeru)}]{\ja{にげる} (nigeru) — to run away; to flee}

\label{lesson:JA-C262-nigeru}



\begin{quote}
Before the new one: say the Japanese for an accident, then the Japanese for to be saved.
\end{quote}

\subsection*{You'll want to know: \ja{にげる}}
\textbf{\ja{にげる}} — \emph{nigeru} — \textquotedblleft{}to run away; to flee\textquotedblright{}.

\begin{quote}
\textbf{\ja{いぬが} \ja{きて、ねこが} \ja{にげました。}} — \emph{inu ga kite, neko ga nigemashita} — \textquotedblleft{}A dog came along, and the cat fled.\textquotedblright{}
\end{quote}

Polite \textbf{\ja{にげます}}, \emph{nigemasu}; an ichidan verb. Past (polite): \textbf{\ja{ねこが} \ja{にげました。}}, \emph{neko ga nigemashita}.

\begin{cousinweb}
\textbf{\ja{にげる}} (\emph{nigeru}) is a native Japanese word (wago), inherited rather than borrowed from Chinese; its kanji was attached later as a kun'yomi (JMdict).
\end{cousinweb}

\begin{grammarlens}[title={What you've built}]
One more word for telling what happened.
\end{grammarlens}

\subsection*{Your turn}
\begin{itemize}
\raggedright
  \item \emph{Say it:} \emph{nigeru}
  \item \emph{Say it:} \emph{nigeru}, once more
  \item \emph{Say it:} \emph{inu ga kite, neko ga nigemashita} — A dog came along, and the cat fled.
  \item \emph{Recall:} say the Japanese for an accident, then the Japanese for to be saved, then say \emph{nigeru} again
\end{itemize}

\begin{tcolorbox}[breakable,colback=teal!4,colframe=teal!35!black,title={Before you move on}]
What does \ja{にげる} mean? (\textquotedblleft{}To run away; to flee\textquotedblright{}.) Say it once more.
\end{tcolorbox}
\section[\texorpdfstring{\ja{とつぜん} (totsuzen)}{とつぜん (totsuzen)}]{\ja{とつぜん} (totsuzen) — suddenly; all of a sudden}

\label{lesson:JA-C262-totsuzen}



\begin{quote}
Before the new one: say the Japanese for to be saved, then the Japanese for to run away.
\end{quote}

\subsection*{You'll want to know: \ja{とつぜん}}
\textbf{\ja{とつぜん}} — \emph{totsuzen} — \textquotedblleft{}suddenly; all of a sudden\textquotedblright{}.

\begin{quote}
\textbf{\ja{とつぜん} \ja{くるまが} \ja{とまりました。}} — \emph{totsuzen kuruma ga tomarimashita} — \textquotedblleft{}All of a sudden the car stopped.\textquotedblright{}
\end{quote}

\begin{cousinweb}
\textbf{\ja{とつぜん}} (\emph{totsuzen}) is Sino-Japanese (kango): two Chinese characters read with their on'yomi, readings borrowed from Chinese. They are \emph{totsu}, \textquotedblleft{}protruding\textquotedblright{} (Middle Chinese \emph{thwot}) and \emph{zen}, \textquotedblleft{}so\textquotedblright{} (\emph{nyen}) (KANJIDIC2; Middle Chinese in Baxter's transcription).
\end{cousinweb}

\begin{grammarlens}[title={What you've built}]
One more word for telling what happened.
\end{grammarlens}

\subsection*{Your turn}
\begin{itemize}
\raggedright
  \item \emph{Say it:} \emph{totsuzen}
  \item \emph{Say it:} \emph{totsuzen}, once more
  \item \emph{Say it:} \emph{totsuzen kuruma ga tomarimashita} — All of a sudden the car stopped.
  \item \emph{Recall:} say the Japanese for to be saved, then the Japanese for to run away, then say \emph{totsuzen} again
\end{itemize}

\begin{tcolorbox}[breakable,colback=teal!4,colframe=teal!35!black,title={Before you move on}]
What does \ja{とつぜん} mean? (\textquotedblleft{}Suddenly; all of a sudden\textquotedblright{}.) Say it once more.
\end{tcolorbox}
\section[\texorpdfstring{\ja{ひがい} (higai)}{ひがい (higai)}]{\ja{ひがい} (higai) — damage; harm suffered}

\label{lesson:JA-C262-higai}



\begin{quote}
Before the new one: say the Japanese for to run away, then the Japanese for suddenly.
\end{quote}

\subsection*{You'll want to know: \ja{ひがい}}
\textbf{\ja{ひがい}} — \emph{higai} — \textquotedblleft{}damage; harm suffered\textquotedblright{}.

\begin{quote}
\textbf{\ja{ひがいは} \ja{ばしょに} \ja{よります。}} — \emph{higai wa basho ni yorimasu} — \textquotedblleft{}How much harm is done depends on the place.\textquotedblright{}
\end{quote}

\begin{cousinweb}
\textbf{\ja{ひがい}} (\emph{higai}) is Sino-Japanese (kango): two Chinese characters read with their on'yomi, readings borrowed from Chinese. They are \emph{hi}, \textquotedblleft{}incur\textquotedblright{} (Middle Chinese \emph{bjeX}) and \emph{gai}, \textquotedblleft{}harm\textquotedblright{} (\emph{hajH}) (KANJIDIC2; Middle Chinese in Baxter's transcription).
\end{cousinweb}

\begin{grammarlens}[title={What you've built}]
One more word for telling what happened.
\end{grammarlens}

\subsection*{Your turn}
\begin{itemize}
\raggedright
  \item \emph{Say it:} \emph{higai}
  \item \emph{Say it:} \emph{higai}, once more
  \item \emph{Say it:} \emph{higai wa basho ni yorimasu} — How much harm is done depends on the place.
  \item \emph{Recall:} say the Japanese for to run away, then the Japanese for suddenly, then say \emph{higai} again
\end{itemize}

\begin{tcolorbox}[breakable,colback=teal!4,colframe=teal!35!black,title={Before you move on}]
What does \ja{ひがい} mean? (\textquotedblleft{}Damage; harm suffered\textquotedblright{}.) Say it once more.
\end{tcolorbox}
